\documentclass[12pt]{article}
\usepackage[margin=1in]{geometry}
\usepackage{parskip}
\usepackage{titlesec}
\usepackage{enumitem}
\usepackage[most]{tcolorbox}
\tcbuselibrary{breakable}
\tcbset{breakable}
\usepackage{hyperref}
\usepackage{bookmark}
\usepackage{setspace}
\usepackage{graphicx}
\usepackage{xcolor}

% Define colors for mechanics
\definecolor{mechblue}{RGB}{0,100,200}
\definecolor{mechgreen}{RGB}{0,150,0}
\definecolor{mechred}{RGB}{200,0,0}
\definecolor{mechgray}{RGB}{80,80,80}

% Mechanics box style
\newtcolorbox{mechanic}[1][]{
  colback=mechgray!5,
  colframe=mechgray!70!black,
  title=#1,
  fonttitle=\bfseries,
  left=2mm,
  right=2mm,
  top=2mm,
  bottom=2mm
}

% Quote style
\renewcommand{\quote}{
  \list{}{%
    \leftmargin=2em%
    \rightmargin=2em%
  }%
  \item[]%
}
\let\endquote=\endlist

% Title formatting
\titleformat{\section}
  {\normalfont\Large\bfseries}{\thesection}{1em}{}
\titleformat{\subsection}
  {\normalfont\large\bfseries}{\thesubsection}{1em}{}
\titleformat{\subsubsection}
  {\normalfont\normalsize\bfseries}{\thesubsubsection}{1em}{}

\begin{document}

\thispagestyle{empty}

\begin{center}
    {\LARGE \bfseries The Mistlands: Bells, Salt, and the Price of Breath}\\[0.5em]
    {\large An Expansion for Fate's Edge}\\[1em]
    {\itshape ``The bells do not mark time. They mark boundaries. When they fall silent, the fog remembers your name.''}\\[2em]
    {\large Revised for Consistency with \textit{Into the Direwood}}\\[0.5em]
    {\normalsize (A Tier III--IV Folk Horror Saga)}\\[1em]
\end{center}

\section*{Introduction: The Land That Was Not Always Mist}
The Mistlands are not a natural place. They were made.

Before the fog, before the bells, before the Aeler came with their keystone contracts and iron wards, this land was called the \textbf{Alder March} -- a fertile, mist‑kissed frontier where Ostrikari herders grazed their horses and Linn raiders came ashore to trade. The grass grew tall enough to hide a child. The rivers ran clear. The dead stayed buried.

Then the warden woke.

His name has been struck from every record, but his crime is remembered: when an Ecktorian legion cut down his sacred grove to build a signal tower, he did not curse the legion. He befriended an exiled necromancer, and together they twisted root and bone into death‑magic. He raised an undead host from every barrow, every battlefield, every forgotten grave. For seven years, the Alder March burned. The Linn heroine -- whose name is now a silence -- drove a gilded spear of unknown origin through his heart on the \textbf{Alder Throne}, a hill of black oak and bone.

But the warden did not die clean.

From his corpse poured the \textbf{Grey Mists} -- a necrotic fog that did not lift with the sun. It crept across the valleys, choking the living, raising the dead as shambling husks. The survivors fled to the coast. The Ostrikari called it the \textit{Wyrd-Weather}; the Linn called it the \textit{Breath of the Hungry}.

Desperate, the chieftains of the Mistlands turned to the only power that could build walls against an immaterial enemy: the \textbf{Aeler}.

The Aeler came. They built the bell‑lines. They blessed the salt pans. They taught the Mistlanders how to ring the wards and keep the fog at bay. And they charged a price -- not in gold, but in \textbf{breath}. A share of every crop, a tithe of every able‑bodied youth sent to work the deep vaults, a perpetual right of audit over every threshold.

The Mistlands have been a protectorate ever since. The people call themselves \textbf{Mistlanders} now -- a blend of Ostrikari, Linn, and other refugees. They keep the bells. They pay the tithe. And they remember, in whispers, the name of the heroine who saved them... and the warden whose curse they still breathe every winter.

\section{What the Mistlands Are Now}
\begin{itemize}[leftmargin=*]
  \item \textbf{Genre:} Folk horror / gothic / supernatural dread
  \item \textbf{Key themes:} Isolation, collective obligation, the cost of safety, forgotten history, the power of remembrance
  \item \textbf{Dominant powers:} Aeler Protectorate (temporal), local bell‑wardens (spiritual), the Direwood (existential), the Forgotten Name (metaphysical)
\end{itemize}

\subsection{The Three Truths of the Mistlands}
\begin{enumerate}[label=\arabic*.]
  \item \textbf{The fog is not weather. It is a wound.} It rises from the Alder Throne, where the warden fell, and it remembers what was forgotten -- specifically, the name of the heroine who bound him.
  \item \textbf{Bells are law.} Every village has a bell‑warden. Every crossing has a bell‑chime. Strike the wrong note and the fog takes notice -- but ring true in remembrance, and the fog hesitates.
  \item \textbf{The Aeler own your breath.} Their wards keep the dead down, but the price is perpetual -- in labor, in loyalty, in the slow erosion of self. Yet their true purpose is not oppression, but preservation: they maintain the wards that hold back the fog \textit{until the heroine's name is remembered}.
\end{enumerate}

\section{History of the Mistlands (GM‑Facing)}

\subsection{The Alder Throne and the Fortress Heart}
A low hill crowned with petrified oak and bleached bone -- but this is only the surface. Beneath, the warden's corpse feeds a forest grown from bone and root: the \textbf{Fortress Heart}. Here, halls pulse faintly as if remembering breath, and at the very center stands the \textbf{stone spear}, planted by the Linn heroine. No one has been able to pull it free. It is said that only one who speaks her true name can draw the spear, and that doing so will end the mists -- or release something worse.

At the foot of the throne stands the spear, planted by the Linn heroine. No one has been able to pull it free. It is said that only one who speaks her true name can draw the spear, and that doing so will end the mists -- or release something worse.

\subsection{The Aeler Compact (Keystone of Suffering)}
Signed in blood and bell‑bronze, this contract established the \textbf{Mistlands Protectorate}. Key clauses:
\begin{itemize}[leftmargin=*]
  \item \textbf{The Tithe of Breath:} Each village must send one young adult per decade to serve in the Aeler deep vaults. Some return; most do not.
  \item \textbf{The Ward‑Salt Levy:} Peat and salt must be given to the Aeler at fixed prices, regardless of harvest.
  \item \textbf{The Right of Audit:} Aeler "Spirit‑Shield Correctors" may enter any threshold at any time to inspect ward integrity. They answer only to the Vault‑Queen.
\end{itemize}

Violating a clause does not bring punishment -- it brings \textbf{ward withdrawal}. The Aeler simply stop maintaining the bell‑lines in that village. The fog does the rest -- but only if the heroine's name remains forgotten. Should her name be spoken in true remembrance, the wards strengthen even without Aeler aid.

\section{Peoples of the Mistlands}

\subsection{Mistlanders (Mixed Ostrikari and Linn)}
The majority population. They speak a creole of Ostrikari breath‑counting and Linn salt‑oaths. Their houses are built with westward-facing doors (windward; the fog comes from the east) and every room has a \textbf{silence bell} -- a small bronze bell that rings if a spirit enters.

\textbf{Cultural touchstones:}
\begin{itemize}[leftmargin=*]
  \item \textbf{Thinblood hospitality:} A guest who shares salt and bread gains one night of safety. Extend to a second night, and the house is obligated to defend them as kin.
  \item \textbf{The Ninth Bell Clause:} No bell may be rung nine times in a row. If it rings nine, the ringer must immediately name a dead person aloud. The fog will hear but may be confused -- unless the name spoken is the heroine"s, in which case the fog recoils.
\end{itemize}

\subsection{Aeler Wardens (The Protectorate)}
Dwarven engineers, clerics, and spirit‑shields who maintain the bell‑lines. They live in fortified \textbf{vent‑keeps} along the eastern ridge. They do not mix with locals. They wear hoods lined with silver mesh to prevent possession.

\textbf{Typical Aeler Warden:}
\begin{itemize}[leftmargin=*]
  \item Speaks Utaran with a clipped accent.
  \item Carries a null‑bell and a ledger of ward‑integrity scores.
  \item Genuinely believes the Compact is fair; has never seen a village fail to pay the Tithe of Breath -- but knows that should the heroine's name be remembered, the tithe may one day be lifted.
\end{itemize}

\subsection{The Unshrouded}
Mistlanders who have rejected the Compact. They live in hidden valleys, trading with smugglers and paying no tithe. They are hunted by both the Aeler and the fog. Some have learned to \textbf{bind small spirits} into iron cages -- a dangerous and heretical art, for binding spirits without remembrance only feeds the fog's hunger for forgotten names.

\subsection{The Linn Heroine's Descendants (If Any)}
Spiral, fractured families who claim descent from the forgotten heroine. They carry a \textbf{white braid} in their hair -- the colour of lost memory. They are respected and pitied. Many become bell‑wardens, hoping to earn enough grace to learn her name -- for they know that only by speaking it can the mist be turned back.

\section{The Direwood (The Wound in the World)}
The Direwood is not a forest. It is the \textbf{warden's curse made forest}. Trees grow from bone. Rivers run grey. Animals have too many eyes. The Direwood covers the eastern half of the Mistlands, and it is expanding. The bell‑lines hold it at bay -- for now, but only because they are maintained in remembrance of what was lost.

\textbf{What lurks within:}
\begin{itemize}[leftmargin=*]
  \item \textbf{Wraiths} -- spirits of the warden's undead host, still marching.
  \item \textbf{The Unsilenced} -- corpses that rise when a bell‑line fails, chanting a hymn in a language no one speaks -- the language of the forgotten name.
  \item \textbf{The Alder King's Vestige} -- a fragment of the corrupted warden, trapped in a hollow oak. It cannot act directly, but it whispers to the desperate: \textit{"Forget her name, and I shall rise again."}
\end{itemize}

\section{New Mechanics}

\subsection{Bell‑Ringer (Minor Talent, 2 XP)}
You have trained with the village bells. Once per scene, when you hear a bell (including your own), you may:
\begin{itemize}[leftmargin=*]
  \item \textbf{Identify the purpose} of the bell‑ring (alarm, warning, celebration, ward‑test) without rolling.
  \item \textbf{Ring a bell} to create a \textbf{Warded [4]} timer in one zone: spirits suffer Position --1 and must test Resolve (DV 3) to enter.
\end{itemize}
\downside: If you ring a bell without first checking the ward‑pattern (i.e., without invoking remembrance), you may tick the \textbf{Fog Attention timer} (see below).

\subsection{Ward‑Salt (Minor Asset, 4 XP)}
A pouch of consecrated salt from the Aeler pan‑works. Use a handful:
\begin{itemize}[leftmargin=*]
  \item \textbf{As a reaction} to a spirit's attack, reduce Harm from that attack by 1 (to a minimum of 0).
  \item \textbf{As an action}, create a \textbf{Salt Line}. A spirit must test Body + Resolve (DV 4) to cross; on a failure, it cannot cross for the scene.
\end{itemize}
\depletion: A pouch has 3 uses. Refill requires a journey to an Aeler vent‑keep and paying the tithe (a memory, a lock of hair, or 1 XP).

\subsection{Fog Attention Clock}
Every time the party fails a roll involving the fog, a spirit, or the Direwood \textit{while forgetting what must be remembered}, the GM may tick the \textbf{Fog Attention [6]} timer. When it fills:
\begin{itemize}[leftmargin=*]
  \item The party is \textbf{Noticed}. All spirit encounters gain +1 Cap, and the fog may manifest a \textbf{Devourer} -- a composite spirit made of three wraiths (Cap 4, Fear aura).
\end{itemize}
Reducing the timer requires a ritual at a bell‑line \textit{while speaking a fragment of the heroine's name} (spend 1 Boon per segment, or a service to the Aeler framed as an act of remembrance).

\subsection{The Tithe of Breath (Narrative Obligation)}
Each Mistlander character (or character who has lived in the Protectorate for more than a year) owes a \textbf{Tithe Clock [4]} to the Aeler. Advance it when:
\begin{itemize}[leftmargin=*]
  \item They use Aeler infrastructure (bell‑lines, ward‑salt, vent‑keep hospitality).
  \item They refuse a direct request from a Spirit‑Shield Corrector \textit{without offering an act of remembrance in lieu}.
\end{itemize}
When the timer fills, they must choose:
\begin{itemize}[leftmargin=*]
  \item \textbf{Serve:} Go to the deep vaults for one month (off‑screen). Clear the timer and gain +2 XP, but return with a \textbf{Shadow} (a minor Condition: --1 die to resist fear, as you"ve seen something that still watches you -- a shadow of the forgotten).
  \item \textbf{Refuse:} The Aeler suspend your access to wards. The character takes an additional \textbf{Fog Attention +2} and becomes a target for Correctors -- unless they perform a public act of remembrance (e.g., speaking a fragment of the heroine's name at a bell‑tower), which may restore partial access.
\end{itemize}

\section{Adventure Seeds (Name Remembrance)}

\subsection{1. The Ninth Bell}
A village bell has been ringing nine times at midnight. The warden is dead. The party must enter the bell‑tower during the ninth ring and silence the spirit that now rings it -- or negotiate a new toll. \textit{But if they can speak a true fragment of the heroine's name while ringing it nine times, the fog will part and reveal a hidden path to the Alder Throne.}

\subsection{2. The Salt Tax}
An Aeler vent‑prior demands double the usual salt quota. The village cannot pay. The prior offers a deal: find the lost \textbf{Spear of the Linn Heroine} in the Direwood, and he will forgive the debt. (He wants to sell it to a collector in Ecktoria.) \textit{However, if the party returns with only the spear shaft -- missing the gilded tip -- the prior will refuse, for the spear's power lies in its wholeness, a symbol of the heroine's unbroken name.}

\subsection{3. The White Braid}
A family claiming descent from the heroine is being hunted by Unshrouded. They believe the warden's vestige has learned a fragment of her true name and wants to erase it forever. Protect them -- or help the Unshrouded steal the braid for their own ritual. \textit{If the braid is destroyed, the heroine's name becomes harder to remember; if it is preserved, the family gains insight into one true syllable.}

\subsection{4. The Forgotten Name (Campaign Arc)}
A Tulkani storyteller arrives with a fragment of a song -- the only surviving syllable of the heroine's name. The party must travel to three locations: the \textbf{Alder Throne}, the \textbf{Sunken Barrow} (where her body was laid), and the \textbf{Echoing Caves} (where her spirit was reflected in a sacred pool). Each location gives a clue. Assembling the name allows one of three choices:
\begin{enumerate}[label=\arabic*.]
  \item \textbf{Speak the name at the Alder Throne} -- the mist ends, but the warden's vestige awakens (new campaign threat).
  \item \textbf{Burn the name} -- the mists remain, but the heroine's spirit is finally released. Gain a \textbf{Blessing of the Lost} (+1 die to all rolls involving fog or spirits for the rest of the campaign).
  \item \textbf{Weave the name into a bell} -- create a new ward‑bell that halves the Aeler's tithe for a generation. The party is hailed as heroes -- but the Aeler take note, for they see their tithe threatened by true remembrance.
\end{enumerate}

\section{GM Toolkit: Running the Mistlands (Updated)}

\subsection{Clocks}
\begin{table}[h]
  \centering
  \begin{tabular}{|l|c|l|l|}
    \hline
    \textbf{Clock} & \textbf{Segments} & \textbf{Advances On} & \textbf{When Full} \\
    \hline
    Fog Attention & 6 & Failed rolls involving fog/spirits \textit{while forgetting remembrance}; ignoring bell‑lines; using magic without salt or remembrance. & A Devourer appears; all spirit caps +1. \\
    \hline
    Aeler Patience & 8 & Refusing audits; smuggling Undocumented goods; sheltering Unshrouded \textit{without acts of remembrance}. & The protectorate sends a \textbf{Corrector} (Cap 5, with a null‑bell and a writ of seizure). \\
    \hline
    The Tithe of Breath (per PC) & 4 & Using Aeler infrastructure; refusing requests \textit{without remembrance}. & Serve or refuse (see above). \\
    \hline
  \end{tabular}
\end{table}

\subsection{Random Mistlands Events (d6)}
\begin{enumerate}[label=\arabic*.]
  \item \textbf{Bell‑line failure:} A segment of the line goes silent. Party must find the broken bell and ring it \textit{with a cadence of remembrance} (Wits + Lore DV 3) or the Fog Attention ticks +1.
  \item \textbf{Wraith procession:} A line of undead marches past, ignoring the living. If the party hides, nothing happens. If they interfere, fight (Cap 3).
  \item \textbf{Salt‑pan fire:} The Aeler's pan‑works are burning. Help put it out (Body + Athletics) to gain a temporary Ward‑Salt use. Fail and the Aeler Patience timer ticks +1.
  \item \textbf{The White Braid's request:} A descendant of the heroine asks the party to escort them to the Alder Throne. No reward, but a future favor \textit{and a fragment of her name revealed in dream}.
  \item \textbf{Fog‑eater:} A creature that breathes mist attacks. It cannot be harmed by steel; only salt, fire, \textit{or a properly rung bell while speaking a name fragment}.
  \item \textbf{The ninth cup:} A tavern keeper offers a ninth cup free. If a PC accepts, the Fog Attention timer ticks +1, \textit{and that night they dream of the heroine's face -- if they speak her name in the dream, they wake with +1 Boon}.
\end{enumerate}

\subsection{The Spear of the Lost Heroine (Epic Asset)}
If the party reclaims the spear from the Alder Throne \textit{requires the heroine's true name}, it gains the following properties:
\begin{itemize}[leftmargin=*]
  \item \textbf{Thrown:} Ranged (Near) attack, Harm 3, ignores spirit armor.
  \item \textbf{Ward‑Breaker:} Once per session, can shatter a [WARD] effect by striking it.
  \item \textbf{Curse:} The wielder is marked by the warden's vestige. Once per session, the GM may tick the Fog Attention timer by 1 at no cost \textit{-- unless the wielder speaks the heroine's name aloud when activating this curse, in which case the vestige recoils and gains no Fog Attention}.
\end{itemize}

\section{Conclusion: The Price of Breath}
The Mistlands are a place of bitter bargains. The fog is not evil -- it is a consequence, still unfolding after centuries. The Aeler are not tyrants -- they are creditors, and the debt is real. The people are not victims -- they are survivors who have learned to count their breaths, salt their thresholds, and ring their bells in just the right cadence.

But the fog is patient. The warden's name is almost forgotten. And somewhere, in the echo of a bell that rang nine times, a Linn heroine waits to be remembered.

\textit{Do you dare speak her name?}

\end{document}
